\chapter{加密资产监管与政策：法域逻辑与执法难点}

\textbf{本章不构成法律意见。}加密规则演进极快，正式展业须以各辖区\textbf{成文法、监管规则与律师意见}为准，并标注检索日期。

\section{欧盟 MiCA：统一框架与稳定币分层}

欧盟《加密资产市场监管条例》（Markets in Crypto-Assets, \textbf{MiCA}，条例编号如 2023/1114）将加密资产定义为可用分布式账本或类似技术以电子方式转移或存储的价值或权利的数字表示，并针对\textbf{其他金融服务立法未涵盖}的加密资产与服务建立统一规则。欧盟委员会概述与立法时间线见：\url{https://finance.ec.europa.eu/digital-finance/crypto-assets_en}；全文见 EUR-Lex：\url{https://eur-lex.europa.eu/legal-content/EN/ALL/?uri=CELEX:32023R1114}。

监管逻辑上，MiCA 对\textbf{资产参考代币（ART）}与\textbf{电子货币代币（EMT）}类稳定币发行人设定更严格的授权、储备、赎回与治理义务；对\textbf{其他加密资产}的发行与营销设定信息披露与白皮书制度；对\textbf{加密资产服务提供商（CASP）}（交易、托管、撮合、顾问等）设定许可与运营要求，并由 ESMA 与各国主管机关协调监督。MiCA 与支付服务、反洗钱（AML）及\textbf{DORA}（数字运营韧性）等规则在机构层面交叉——持牌金融机构参与链上业务时，常需同时满足多条欧盟立法。

\section{美国：证券与商品的「双头监管」叙事}

在美国，同一数字资产是否构成\textbf{证券}（归 SEC 管辖的投资合同分析，常引用 \textbf{Howey} 判例标准：金钱投资、共同事业、对他人努力的合理期待利润）与是否属于\textbf{商品}（CFTC 对特定虚拟货币的界定与衍生品监管）长期存在事实与法律争议。\textbf{现货}加密交易平台、经纪与托管是否需注册为经纪交易商、交易所或适用州级\textbf{货币传输（money transmission）}许可，取决于具体商业模式与资产分类。近年立法与执法动向频繁（包括稳定币专项立法讨论、ETF 批准条件、执法和解与规则提案），本书无法逐条固化；读者应订阅 SEC、CFTC、FinCEN、OCC、美联储与州监管机构更新，并由美国证券律师出具产品备忘录。

\section{稳定币与支付稳定币：政策焦点}

稳定币连接\textbf{链上 DeFi}与\textbf{链下法币体系}，政策关切集中在：储备资产质量与透明度、赎回压力下的流动性、发行人破产时持有人顺位、系统性重要性（多链与多应用共用同一稳定币）、以及 AML/制裁合规。欧盟 MiCA 对 ART/EMT 的强化规则即体现该思路。美国层面是否存在联邦统一框架及与州法的优先关系，以\textbf{当时已签署生效的联邦法与实施细则}为准——设计产品者须做「双轨」合规评估而非依赖单一博客解读。

\section{FATF 与「旅行规则」}

金融行动特别工作组（FATF）对虚拟资产服务提供商（VASP）提出与银行类似的\textbf{AML/CFT} 期望，包括\textbf{旅行规则（Travel Rule）}：在转账超过阈值时交换并记录发起方与受益方信息。各国转化程度不同，但机构级托管、交易所与链上分析工具已广泛部署交易监控与地址标注。DeFi 无许可接口与自托管钱包使\textbf{合规边界}模糊，成为国际讨论中的难点。

\section{亚太摘录}

\textbf{香港：}虚拟资产交易平台（VATP）等事项由证监会（SFC）等机关规管；规则与持牌名单以官网为准：\url{https://www.sfc.hk/en/Rules-and-standards/Virtual-assets/Virtual-asset-trading-platforms-operators}。实务上常同时涉及证券型代币与非证券型代币业务线条，申请与持续合规成本较高。

\textbf{新加坡：}支付型数字代币服务可能落入\textbf{PSA} 下牌照与豁免框架；证券型代币仍适用《证券期货法》。MAS 对零售加密营销与稳定币监管持续收紧，须查阅现行指南。

\textbf{中国大陆：}对虚拟货币交易炒作、代币发行融资等采取严格限制立场；境内机构与个人参与境外平台与代币发行仍可能涉及多重法律风险。本书读者若以人民币区为主要基地，须单独取得境内律师书面意见。

\section{DeFi、DAO 与「无许可」执法难题}

协议可匿名部署，治理代币可全球分发，物理运营团队可分散司法辖区——监管机构可能选择\textbf{对可识别主体}（前端运营者、核心开发者、营销方、稳定币发行人、托管入口）执法，或通过\textbf{制裁与域名/基础设施封锁}施加压力。产品设计若假设「链上无监管」，常与\textbf{入口合规}（法币出入金、应用商店、API 供应商）冲突。

\begin{figure}[htbp]
\centering
\begin{tikzpicture}[font=\scriptsize]
  \node[ellipse, draw=blue!70, fill=blue!5, minimum width=3.6cm, minimum height=2.2cm, align=center] (eu) {\textbf{欧盟}\\MiCA/CASP\\ART/EMT};
  \node[ellipse, draw=teal!70, fill=teal!5, minimum width=3.6cm, minimum height=2.2cm, align=center, right=1.8cm of eu] (us) {\textbf{美国}\\SEC/CFTC\\州法/FinCEN};
  \path (eu) -- node[coordinate, midway] (midewus) {} (us);
  \node[ellipse, draw=orange!70, fill=orange!5, minimum width=3.6cm, minimum height=2.2cm, align=center, below=1.1cm of midewus] (ap) {\textbf{亚太}\\MAS/SFC 等\\VASP/PSA};
  \draw[compliance arrow, dashed] (eu) -- (us);
  \draw[compliance arrow, dashed] (eu.south) -- (ap.north west);
  \draw[compliance arrow, dashed] (us.south) -- (ap.north east);
  \node[below=0.15cm of ap, text width=8.5cm, align=center, text=gray!55!black]
    {跨境业务常触发\textbf{多法域并行适用}；代币「标签」不改变经济实质认定。};
\end{tikzpicture}
\caption{监管地理示意（高度简化；非穷举）}
\label{fig:reg-geo-web3}
\end{figure}

\section{与私募、资管业务的接口}

若基金或管理人计划：配置加密现货/衍生品、投资代币化基金份额、运行节点或质押策略、或向 LP 提供链上收益产品，须在 LPA 与侧信中披露\textbf{策略边界与风险}，并完成本书合规篇所述\textbf{管理人资格与产品注册/备案}在加密场景下的增量论证。代币化证券若属于「证券」，通常须遵循传统证券发行与交易平台规则，而非仅依赖「链上自治」叙事。
